%!Mode:: "TeX:UTF-8"
%!TEX encoding = UTF-8 Unicode
%arara: xelatex
\documentclass{ctexart}
\newif\ifpreface
%\prefacetrue
\input{../../../global/all}
\begin{document}
\large
\setlength{\baselineskip}{1.2em}
\ifpreface
    \input{../../../global/preface}
\newgeometry{left=2cm,right=2cm,top=2cm,bottom=2cm}
\else
\newgeometry{left=2cm,right=2cm,top=2cm,bottom=2cm}
\maketitle
\fi
%from_here_to_type
\begin{problem} 
从 $k$ 步 Krylov--Schur 分解到 $m$ 步 Arnoldi 过程.
\end{problem}
\begin{solution}
设 $A\in\mathbb{R}^{n\times n}$。
$k$ 步 Krylov--Schur 分解具有如下形式：
\begin{equation}
A U_k = U_k H_k + u_{k+1} h_{k+1,k} e_k^{\top},
\label{eq:ks}
\end{equation}
其中
$U_k=[u_1,\dots,u_k]\in\mathbb{R}^{n\times k}$ 为列正交矩阵，
$H_k\in\mathbb{R}^{k\times k}$ 为上 Hessenberg 矩阵，
$u_{k+1}\perp U_k$ 且 $\|u_{k+1}\|_2=1$，
$e_k$ 为第 $k$ 个标准基向量。

等价地，式 \eqref{eq:ks} 可写为块矩阵形式
\begin{equation}
A U_k
=
\begin{bmatrix}
U_k & u_{k+1}
\end{bmatrix}
\begin{bmatrix}
H_k \\[1mm]
h_{k+1,k} e_k^{\top}
\end{bmatrix}.
\end{equation}

在此基础上，对向量 $u_{k+1}$ 继续执行 Arnoldi 正交化过程，
生成正交向量 $u_{k+2},\dots,u_{m+1}$，
从而得到 $m$ 步 Arnoldi 分解。
记
\[
U_m = [u_1,\dots,u_m]\in\mathbb{R}^{n\times m},
\qquad
U_{m+1}=[U_m,\;u_{m+1}],
\]
则有
\begin{equation}
A U_m = U_{m+1}\,\bar H_m,
\label{eq:arnoldi}
\end{equation}
其中 $\bar H_m\in\mathbb{R}^{(m+1)\times m}$ 为上 Hessenberg 矩阵。

由于前 $k$ 个基向量来源于 Krylov--Schur 分解，
矩阵 $\bar H_m$ 具有如下块结构：
\begin{equation}
\bar H_m =
\begin{bmatrix}
H_k & * \\
0   & H_{22} \\
0   & h_{m+1,m} e_{m-k}^{\top}
\end{bmatrix},
\end{equation}
其中 $H_{22}\in\mathbb{R}^{(m-k)\times(m-k)}$ 为后续 Arnoldi 过程产生的
上 Hessenberg 矩阵，符号 $*$ 表示子空间之间的耦合项。

因此，从 $k$ 步 Krylov--Schur 分解扩展到 $m$ 步 Arnoldi 过程后，
得到的整体矩阵形式为
\begin{equation}
A U_m
=
\begin{bmatrix}
U_k & U_{k+1:m} & u_{m+1}
\end{bmatrix}
\begin{bmatrix}
H_k & * \\
0   & H_{22} \\
0   & h_{m+1,m} e_{m-k}^{\top}
\end{bmatrix}.
\end{equation}
\end{solution}



\end{document}
